
\begin{titlepage}
\thispagestyle{empty}

\centering

\vspace{1cm}

\begin{minipage}{0.9\textwidth}
    \hspace*{1em}
    코드 정적 분석은 프로그램을 실행하지 않고 코드를 논리적으로 분석하여 동적 특성을 파악하는
기법이며 런타임 오류, 주어진 명세 준수 등 다양한 논리적/의미적 성질을 검증하기 위해 사용된다.
\traceinspector는 요약해석(abstract interpretation)이론을 바탕으로 구현된 추적 가능한 시각화 기반 코드 정적분석기이다.
단순히 오류를 탐지해서 알려주는 정적 분석기를 넘어, 분석 과정을 시각적으로
나타내어 사용자가 어떤 과정을 통해 오류가 탐지되었는지 해석 가능하게끔 돕고자 한다.
\end{minipage} \par

\end{titlepage}